\documentclass[../../main.tex]{subfiles}
\begin{document}

{\bf [解答]} 簡単のため
\begin{align}
  \begin{cases}
    f(x)=x^2-5x \\
    g(x)=\cfrac{\pi}{5}\sin\left(\cfrac{\pi x}{5}\right)-\cfrac{3}{5}\sin^2\left(\cfrac{\pi x}{5}\right)
  \end{cases}
\end{align}
とおく．$0<x<5$では$f(x)<0$，$g(x)>0$である．
故にグラフの概形は下図である．

\begin{figure}[htb]
  \centering
  \begin{tikzpicture}[x=1.0cm, y=0.52cm, >=stealth, every node/.style={font=\small}]
    \coordinate (O) at (0, 0);
    \coordinate (A) at (5, 0);
    \coordinate (V) at (2.5, -6.25);

    % x軸上側の g(x)>0 の概形・領域 S_2 が判読できるよう g(x) の縦方向のみ5倍に拡大して図示
    % S_2 (上側領域: 0 < y < g(x))
    \fill[pattern=north west lines, pattern color=teal!75!black]
    plot[domain=0:5, samples=80, smooth] (\x, {5*(0.6283185*sin(36*\x) - 0.6*pow(sin(36*\x), 2))})
    -- (5, 0) -- (0, 0) -- cycle;

    % S_1 (下側領域: f(x) < y < 0)
    \fill[pattern=north east lines, pattern color=blue!55]
    (0, 0) -- plot[domain=0:5, samples=80, smooth] (\x, {\x*\x - 5*\x}) -- cycle;

    % 座標軸
    \draw[->, thick, gray!80!black] (-0.7, 0) -- (5.9, 0) node[right] {$x$};
    \draw[->, thick, gray!80!black] (0, -6.9) -- (0, 1.8) node[above] {$y$};
    \node[below left=2pt] at (O) {$O$};

    % 頂点からの点線垂線
    \draw[densely dotted, thick, gray!70!black] (2.5, 0) -- (V) -- (0, -6.25);
    \node[above=1pt, font=\scriptsize, fill=white, inner sep=1pt] at (2.5, 0) {$\frac{5}{2}$};
    \node[left=2pt, font=\scriptsize] at (0, -6.25) {$-\frac{25}{4}$};

    % 曲線 y = f(x) と y = g(x)
    \draw[line width=1.4pt, blue!85!black, domain=-0.3:5.3, samples=90, smooth]
    plot (\x, {\x*\x - 5*\x}) node[right, font=\footnotesize] {$y=f(x)$};

    \draw[line width=1.3pt, teal!85!black, domain=-0.25:5.25, samples=100, smooth]
    plot (\x, {5*(0.6283185*sin(36*\x) - 0.6*pow(sin(36*\x), 2))});
    \node[above right=1pt, font=\footnotesize, teal!85!black] at (0.85, 0.82) {$y=g(x)$};

    % 面積ラベルと端点
    \node[font=\footnotesize, blue!85!black, fill=white, inner sep=1.5pt] at (2.5, -2.9) {$S_1$};
    \draw[<-, thick, teal!85!black] (4.15, 0.55) to[bend left=18] (4.7, 1.35)
    node[right, font=\footnotesize, teal!85!black] {$S_2$};

    \node[circle, fill=gray!80!black, inner sep=1.3pt] at (O) {};
    \node[circle, fill=gray!80!black, inner sep=1.3pt] at (A) {};
    \node[above right=1pt, font=\scriptsize] at (A) {$5$};
    \node[circle, fill=blue!85!black, inner sep=1.3pt] at (V) {};
  \end{tikzpicture}
  \caption{$y=f(x)$ と $y=g(x)$ が囲む領域（$S=S_1+S_2$）}
  \label{fig:region_f_and_g}
\end{figure}
これら二つのグラフの囲む面積$S$とすると，
\begin{align}
  S=\int_0^5(g(x)-f(x))dx\label{eq:1}
\end{align}
である．各項計算すると，
\begin{align}
  &\cdot S_1=-\int_0^5f(x)dx=\dfrac{5^3}{6}=\dfrac{125}{6} \\
  &\cdot S_2=\int_0^5g(x)dx=\int_0^\pi\left(\sin t-\dfrac{3}{\pi}\sin^2t\right)dt \\
  &=\left[-\cos t-\dfrac{3}{2\pi}\left(t-\dfrac{1}{2}\sin 2t\right)\right]_0^{\pi}
  =2-\dfrac{3}{2}=\dfrac{1}{2}
\end{align}
である．故に\cref{eq:1}に代入して
\begin{align}
  S=S_1+S_2=\dfrac{64}{3}\label{eq:2}
\end{align}
である．

また，$S_1>S_2$であるから，題意より，$\alpha<0$となることが必要で,このとき，$y=\alpha x$と$y=f(x)$が$0<x<5$に交点をひとつ持つ．
グラフの概形は下図である．

\begin{figure}[htb]
  \centering
  \begin{tikzpicture}[x=1.0cm, y=0.52cm, >=stealth, every node/.style={font=\small}]
    \coordinate (O) at (0, 0);
    \coordinate (A) at (5, 0);
    \coordinate (P) at (4, -4);

    % 上半分の領域 (y = \alpha x より上側，面積 S/2)
    \fill[pattern=north west lines, pattern color=blue!45]
    plot[domain=0:5, samples=80, smooth] (\x, {5*(0.6283185*sin(36*\x) - 0.6*pow(sin(36*\x), 2))})
    -- plot[domain=5:4, samples=25, smooth] (\x, {\x*\x - 5*\x})
    -- (0, 0) -- cycle;

    % 下半分の領域 (f(x) < y < \alpha x, 0 < x < p，面積 S/2)
    \fill[pattern=north east lines, pattern color=red!75!black]
    (0, 0) -- plot[domain=0:4, samples=65, smooth] (\x, {\x*\x - 5*\x})
    -- (4, -4) -- cycle;

    % 座標軸
    \draw[->, thick, gray!80!black] (-0.7, 0) -- (5.9, 0) node[right] {$x$};
    \draw[->, thick, gray!80!black] (0, -6.9) -- (0, 1.8) node[above] {$y$};
    \node[below left=2pt] at (O) {$O$};

    % 交点 P(p, \alpha p) からの点線垂線
    \draw[densely dotted, thick, gray!70!black] (4, 0) -- (P) -- (0, -4);
    \node[above=1pt, font=\scriptsize, fill=white, inner sep=1pt] at (4, 0) {$p$};
    \node[left=2pt, font=\scriptsize] at (0, -4) {$\alpha p$};

    % 曲線 y = f(x), y = g(x) と直線 y = \alpha x
    \draw[line width=1.4pt, blue!85!black, domain=-0.3:5.3, samples=90, smooth]
    plot (\x, {\x*\x - 5*\x}) node[right, font=\footnotesize] {$y=f(x)$};

    \draw[line width=1.2pt, teal!85!black, domain=-0.25:5.25, samples=100, smooth]
    plot (\x, {5*(0.6283185*sin(36*\x) - 0.6*pow(sin(36*\x), 2))});
    \node[above right=1pt, font=\footnotesize, teal!85!black] at (0.85, 0.82) {$y=g(x)$};

    \draw[line width=1.4pt, red!85!black] (-0.45, 0.45) -- (5.35, -5.35)
    node[right, font=\footnotesize] {$y=\alpha x$};

    % 面積ラベルと交点プロット
    \node[font=\footnotesize, red!85!black, fill=white, inner sep=1.5pt] at (2.1, -3.9) {$\frac{1}{2}S$};
    \node[font=\footnotesize, blue!85!black, fill=white, inner sep=1.5pt] at (2.3, -1.1) {$\frac{1}{2}S$};

    \node[circle, fill=gray!80!black, inner sep=1.3pt] at (O) {};
    \node[circle, fill=gray!80!black, inner sep=1.3pt] at (A) {};
    \node[above right=1pt, font=\scriptsize] at (A) {$5$};
    \node[circle, fill=red!85!black, inner sep=1.5pt] at (P) {};
  \end{tikzpicture}
  \caption{直線 $y=\alpha x$（$\alpha<0$）による領域の二等分と交点 $x=p$}
  \label{fig:bisected_region_by_line}
\end{figure}

交点の$x$座標を$p$とおく．つまり
\begin{align}
  &
  \begin{cases}
    0<p<5 \\
    \alpha p=p^2-5p
  \end{cases}\nonumber \\
  \Longleftrightarrow
  &
  \begin{cases}
    0<p<5 \\
    \alpha =p-5
  \end{cases}\label{eq:3}
\end{align}
とする．\cref{eq:2}に注意して，面積が等しい条件から
\begin{align}
  &\dfrac{1}{2}S=\int_0^p(\alpha x-f(x))dx \\
  \Longleftrightarrow &\dfrac{32}{3}=\dfrac{1}{6}p^3\\
  \Longleftrightarrow &p=4
\end{align}
である．これを\cref{eq:3}に代入して
\begin{align}
  \alpha=-1
\end{align}
となる．
これは$\alpha<0$を満たして十分である．
以上から求める値は$\alpha=-1$である．$\cdots$(答)
\newpage
\end{document}
